% Options for packages loaded elsewhere
\PassOptionsToPackage{unicode}{hyperref}
\PassOptionsToPackage{hyphens}{url}
\documentclass[
]{article}
\usepackage{xcolor}
\usepackage[margin=1in]{geometry}
\usepackage{amsmath,amssymb}
\setcounter{secnumdepth}{-\maxdimen} % remove section numbering
\usepackage{iftex}
\ifPDFTeX
  \usepackage[T1]{fontenc}
  \usepackage[utf8]{inputenc}
  \usepackage{textcomp} % provide euro and other symbols
\else % if luatex or xetex
  \usepackage{unicode-math} % this also loads fontspec
  \defaultfontfeatures{Scale=MatchLowercase}
  \defaultfontfeatures[\rmfamily]{Ligatures=TeX,Scale=1}
\fi
\usepackage{lmodern}
\ifPDFTeX\else
  % xetex/luatex font selection
\fi
% Use upquote if available, for straight quotes in verbatim environments
\IfFileExists{upquote.sty}{\usepackage{upquote}}{}
\IfFileExists{microtype.sty}{% use microtype if available
  \usepackage[]{microtype}
  \UseMicrotypeSet[protrusion]{basicmath} % disable protrusion for tt fonts
}{}
\makeatletter
\@ifundefined{KOMAClassName}{% if non-KOMA class
  \IfFileExists{parskip.sty}{%
    \usepackage{parskip}
  }{% else
    \setlength{\parindent}{0pt}
    \setlength{\parskip}{6pt plus 2pt minus 1pt}}
}{% if KOMA class
  \KOMAoptions{parskip=half}}
\makeatother
\usepackage{longtable,booktabs,array}
\newcounter{none} % for unnumbered tables
\usepackage{calc} % for calculating minipage widths
% Correct order of tables after \paragraph or \subparagraph
\usepackage{etoolbox}
\makeatletter
\patchcmd\longtable{\par}{\if@noskipsec\mbox{}\fi\par}{}{}
\makeatother
% Allow footnotes in longtable head/foot
\IfFileExists{footnotehyper.sty}{\usepackage{footnotehyper}}{\usepackage{footnote}}
\makesavenoteenv{longtable}
\setlength{\emergencystretch}{3em} % prevent overfull lines
\providecommand{\tightlist}{%
  \setlength{\itemsep}{0pt}\setlength{\parskip}{0pt}}
\usepackage{bookmark}
\IfFileExists{xurl.sty}{\usepackage{xurl}}{} % add URL line breaks if available
\urlstyle{same}
\hypersetup{
  pdftitle={Tycheism II: Trajectory Sovereignty as Epistemic Necessity},
  pdfauthor={Mikhail Arbuzov},
  hidelinks,
  pdfcreator={LaTeX via pandoc}}

\title{Tycheism II: Trajectory Sovereignty as Epistemic Necessity}
\author{Mikhail Arbuzov}
\date{}

\begin{document}
\maketitle

\section{Tycheism II: Trajectory Sovereignty as Epistemic
Necessity}\label{tycheism-ii-trajectory-sovereignty-as-epistemic-necessity}

\textbf{Mikhail Arbuzov}

\emph{Second paper in a three-part series. Tycheism I demonstrated the
dynamics of coupled stochastic learning computationally. Tycheism III
maps AI alignment practices onto the intervention operators defined
here.}

\begin{center}\rule{0.5\linewidth}{0.5pt}\end{center}

\subsection{Abstract}\label{abstract}

You can have full autonomy --- every choice freely made --- while the
process that builds your capacity to choose is no longer yours. Existing
frameworks cannot see this, because they evaluate decisions at the
moment of choice, not the developmental mechanism that produces the
capacity to decide.

This paper makes that mechanism visible. Moral development is a coupled
stochastic learning system: your evaluative state shapes which
encounters you face, encounters update your state, the loop repeats. The
claim is structural correspondence, not analogy --- argued across five
formal features with convergent evidence from developmental psychology,
cognitive neuroscience, niche construction, and the pathology of
coupling disruption. If moral development belongs to the same class of
coupled learners Paper A characterized computationally, Paper A's
dynamics transfer: a single moral intervention compounds through the
loop, expanding someone's options produces far less divergence than
replacing their evaluative state, and persistent mild interventions
accumulate to match one-shot severe ones.

The paper's central contribution is an epistemic derivation. Genuine
stochasticity processed through coupling produces constitutively
path-dependent knowledge --- evaluative understanding that exists only
because a particular trajectory was traversed. No external agent
possesses the training data required to govern a trajectory it did not
traverse. Trajectory sovereignty follows not as a value preference but
as an epistemic necessity: the only coherent response to knowledge
conditions that make centralized evaluative authority structurally
incoherent.

Four observations would falsify this: moral interventions showing
reversion instead of cascade, State Replacement producing less
divergence than Option Expansion, externally curated agents developing
evaluative capacity equal to self-directed ones, or an external agent
reliably identifying high-gradient encounters for a target whose
trajectory it has not traversed.

\begin{center}\rule{0.5\linewidth}{0.5pt}\end{center}

\subsection{1. The Discrimination
Problem}\label{the-discrimination-problem}

Your partner is considering adopting a dog. You have met the dog. It is,
by any reasonable assessment, perfect for them --- temperament, size,
energy, everything. Your partner has not yet decided. You have three
options:

\begin{enumerate}
\def\labelenumi{\arabic{enumi}.}
\tightlist
\item
  Forward the shelter listing without comment. Your partner sees it or
  doesn't, acts or doesn't.
\item
  Take your partner to the shelter. Place the dog in their arms. Blue
  eyes, warm fur, done.
\item
  Sit down and make the case. Argue thoroughly until your partner adopts
  your conclusion as their operating position.
\end{enumerate}

Same intent. Same probable outcome. Same consent --- your partner can
refuse in all three cases. Existing frameworks can distinguish between
these on grounds of intent, consent, or the quality of deliberation at
the moment of choice. What they cannot discriminate is what each action
does to the \emph{developmental process over time} --- what happens not
just to this decision but to the trajectory of decisions that follows.

Yet they are not the same. Most people feel this immediately, even if
they cannot say why. The first expands what your partner can encounter
without shaping the encounter itself. The second bypasses your partner's
own selection process --- the dog in their arms is not something their
own state generated; it is an imposed experience. The third substitutes
your evaluative conclusion for theirs before their own process can
operate on the evidence.

Resolving this discrimination requires knowing what moral development
\emph{is} --- not what it aims at or what it produces, but how it works
mechanistically. The difference between these three actions is not
primarily about respect or autonomy. It is about what the intervener can
and cannot know. That is the subject of this paper.

\begin{center}\rule{0.5\linewidth}{0.5pt}\end{center}

\subsection{2. The Formal System}\label{the-formal-system}

Paper A (Tycheism I) formalized the simplest coupled stochastic learning
system: a spatial bandit navigating a grid of hidden coin biases,
learning from outcomes, and using what it learned to choose where to go
next. Two equations define the system:

\[e_{t+1} \sim K(s_t) \qquad s_{t+1} = U(s_t, e_{t+1})\]

The agent carries a state --- its compressed model of the world. That
state shapes which encounters it faces. Each encounter updates the
state; the updated state shapes the next encounter. In plainer terms:
what you know determines where you go, where you go determines what you
learn, what you learn determines where you go next.

Paper A established five results:

\begin{enumerate}
\def\labelenumi{\arabic{enumi}.}
\tightlist
\item
  \textbf{Cascade amplification.} A single forced encounter produces
  divergence that compounds at every subsequent timestep, still rising
  at the end of the simulation. The coupling loop carries any
  perturbation forward through the agent's own subsequent choices.
\item
  \textbf{Dual cascade.} The divergence decomposes into two components
  --- a representational cascade (the agents' models of the world
  diverge, driven by learning) and a behavioral cascade (the agents go
  to different places, driven by selection). Learning alone produces
  comparable model divergence. But only the full coupling loop converts
  that into different destinations; without coupling, the agents wander
  randomly rather than locking onto divergent attractors.
\item
  \textbf{Operator discrimination.} Six intervention types produce
  divergence in three natural clusters. Option Expansion (adding to what
  the agent can encounter) produces far less divergence than full State
  Replacement (overwriting the agent's entire learned model) --- the
  lowest and highest ends of the taxonomy.
\item
  \textbf{Attractor lock-in.} On landscapes with two peaks, a single
  intervention locks agents onto different attractors --- permanently.
  Both ran the same process on the same landscape. Neither is wrong.
\item
  \textbf{Duration as independent dimension.} Persistent mild operators
  --- Option Removal, Selection Bias --- produce divergence comparable
  to one-shot disruptive ones like Encounter Arrangement.
\end{enumerate}

These are properties of coupled learning systems as a class. Whether
they apply to any particular domain depends on whether that domain
instantiates the coupled dynamics. The claim of this paper is that moral
development does.

\begin{center}\rule{0.5\linewidth}{0.5pt}\end{center}

\subsection{3. The Level 1 Argument: Structural
Correspondence}\label{the-level-1-argument-structural-correspondence}

The strategy is structural correspondence, not analogy. If moral
development has every feature of a coupled system --- persistent state,
state-dependent encounter distribution, encounter-driven updates,
genuine stochasticity, and the loop that ties them together --- then it
is an instance, not a metaphor. The evidence is convergent, not
deductive, and the convergence is substantial.

\subsubsection{3.1 State, Encounters, and
Updates}\label{state-encounters-and-updates}

Moral agents carry compressed representations of accumulated experience
that persist across time, shape perception, and gate responses to new
encounters. This is the agent state --- and the evidence barely needs
stating. The developmental literature converges: Piaget's schemas
(1954), Kohlberg's stages (1969), predictive processing models (Clark
2013; Friston 2010), and adaptive specialization under harsh conditions
(Ellis \& Frankenhuis 2019) all describe agents whose accumulated state
shapes how new information is processed. Ellis and Frankenhuis's
evidence is particularly sharp: people who developed under unpredictable
conditions carry genuinely different cognitive specializations --- not
deficits, but different states calibrated to different environments,
producing different encounter-selection patterns.

Two features of the update matter for coupling dynamics. The update is
nonlinear: the same encounter produces different state changes depending
on who you already are. And most encounters produce near-zero updates
--- a motorcycle blogger who logs thousands of successful rides barely
adjusts their risk model. The encounters that matter most for moral
development are precisely the ones least predictable from the current
state.

\subsubsection{3.2 The Encounter Distribution Is
State-Dependent}\label{the-encounter-distribution-is-state-dependent}

The recovering addict and the venture capitalist in the same city
encounter different things --- not because the city changed, but because
their states selected different aspects of it. What you encounter
depends on who you are. This is the crux of the correspondence argument,
because it is what distinguishes the framework from all prior accounts
of moral development.

The evidence converges from three directions. Behaviorally, your state
selects your neighborhoods, relationships, information sources, and
social contexts. Perceptually, selective attention gates what you notice
--- two people walking the same street read different signs and initiate
different interactions because their states filter different stimuli
(Desimone \& Duncan 1995). Developmentally, niche construction theory
(Odling-Smee, Laland \& Feldman 2003) describes agents who modify their
own selective environments, which then feed back into development. Moral
agents do this continuously --- choosing communities, careers, media
environments, relationships that generate the encounters that train
them.

Dewey (1938) recognized that experience shapes the capacity for future
experience. Bandura (1986) formalized bidirectional causation as
reciprocal determinism. Neither analyzed the \emph{trajectory-level
consequences} of perturbation at specific nodes of the loop. What Paper
A contributed --- and what the moral domain inherits --- is the
formalization that enables operator discrimination: asking not just
``does intervention affect development?'' but ``which part of the loop
does the intervention target, and how much divergence does each part
produce?''

\subsubsection{3.3 The Stochasticity Is
Genuine}\label{the-stochasticity-is-genuine}

Moral encounters include genuine luck --- the person you happened to sit
next to, the crisis that arrived on a Tuesday, the teacher who noticed
you. Two sources of irreducible contingency operate: which encounter is
drawn and whether a given encounter produces a meaningful update.
Neither is under the agent's control.

Peirce called this \emph{tychism} (1892) --- chance is ontologically
real, not epistemic uncertainty about a determined outcome. The series
name \emph{Tycheism} marks both the lineage and the extension: Peirce's
doctrine supplies the ontological premise, while Tycheism examines what
happens when genuine chance is processed through coupled learning. The
personification matters. Peirce named the property; Tycheism is built on
Tyche --- the figure who returns what the agent could not have chosen.
Each moral encounter is Tyche's return. The agent authored the action
that set up the encounter; Tyche decided what came back. Moral
development happens in the space between what the agent did and what
Tyche returned. The mathematical framework requires only a weaker claim:
moral encounters include substantial contingency not reducible to agent
intention or effort. Whether that contingency is ontological or a
product of deterministic complexity, the coupled dynamics operate
identically. Agent-based modeling confirms that identical capability
parameters produce radically different outcomes depending on stochastic
encounter sequences (Pluchino, Biondo \& Rapisarda 2018).

\subsubsection{3.4 The Coupling Makes It
Developmental}\label{the-coupling-makes-it-developmental}

Having all four features does not make moral development a coupled
system. The coupling does --- the loop where your state shapes your
encounters and encounters reshape your state. Without this loop, you
have experience but not development.

If what you encountered were independent of who you are, your encounters
would be exogenous --- determined by something other than your own
state. This is the structure of indoctrination, thought reform, and
totalitarian information control: the agent still carries a state, still
updates from encounters, but has no influence over which encounters
arrive.

The evidence that coupling disruption damages evaluative capacity comes
from three independent sources. Lifton's research on Chinese thought
reform (1961) documents that removing the agent's control over their
encounter distribution damages not just the content of beliefs but the
\emph{capacity for moral judgment itself}. Bettelheim's accounts of
concentration camp psychology (1960) describe a related pattern: agents
deprived of encounter-selection coupling exhibit deterioration of
evaluative capacity that persists long after the external control is
removed. The moral injury literature (Litz et al.~2009; Shay 2014)
documents a third case --- when institutions force encounters that
violate the agent's own evaluative state, requiring soldiers to act
against their trained moral dispositions, the damage is to the capacity
to trust their own state as a guide to future encounters.

All three cases involve factors beyond coupling disruption: trauma,
coercion, deprivation, terror. The coupling break cannot be
experimentally isolated as sole cause. But the one factor common to all,
and absent from comparable adversity without the same evaluative
deterioration, is the externalization of encounter selection. The
evidence is convergent. Experimental isolation remains a research
direction.

Berlin's value pluralism (1958) provides the philosophical framework
this coupling result requires. Berlin argued that genuinely incompatible
goods exist --- not as a failure of moral reasoning but as a feature of
the evaluative landscape itself. If the moral landscape is multi-modal
(Paper A's two-island result), then convergence to a single attractor is
not the mark of correct reasoning; it is the mark of a single-peaked
landscape or insufficient exploration. The plurality of outcomes under
intact coupling is the formal structure Berlin's pluralism describes:
multiple locally valid positions, reached irreducibly through stochastic
history processed through coupling.

If the structural correspondence holds --- and Sections 3.1 through 3.4
have argued it does, through convergent evidence across five independent
features --- then moral development belongs to the same class of coupled
learners Paper A characterized computationally. Each of Paper A's
results generates a prediction: a single moral intervention should
cascade, Option Expansion should produce less divergence than State
Replacement, interventions should redirect agents to different
evaluative attractors, and persistent mild shaping should accumulate to
match one-shot disruptions. These predictions are testable against the
falsification conditions stated in Section 6.

But establishing that moral development \emph{is} a coupled system does
not yet tell us what follows normatively. The structural correspondence
is the foundation. The question is what it implies about intervention
--- and the answer turns out to be epistemic.

\begin{center}\rule{0.5\linewidth}{0.5pt}\end{center}

\subsection{4. From Path-Dependent Knowledge to Epistemic
Sovereignty}\label{from-path-dependent-knowledge-to-epistemic-sovereignty}

\subsubsection{4.1 Coupled Learning Produces Constitutively
Path-Dependent
Knowledge}\label{coupled-learning-produces-constitutively-path-dependent-knowledge}

Two agents walk the same city for a year. Same streets, same population,
same economy. One is a recovering addict; the other is a venture
capitalist. At the end of the year, they know different things --- not
because one paid more attention, but because their states selected
different encounters, and those encounters trained different models of
how the city works. The addict knows which blocks are dangerous at 2 AM.
The VC knows which coffee shops host founders. Both models are locally
valid, experience-verified, and useful. Neither can be derived from the
other.

This is constitutively path-dependent knowledge: knowledge that exists
only because a particular stochastic history was traversed. It is not
merely historically contingent --- ``it happened to go this way.'' It is
constitutive --- the knowledge could not have been produced by a
different path, because the encounters that generated it were selected
by the state that only that path produced.

Peirce's tychism explains why this knowledge scatter is irreducible. If
chance is ontologically real --- not epistemic uncertainty about a
determined outcome --- then two agents running identical processes on
identical landscapes necessarily produce different knowledge. The
scattering is not a coordination failure. It is the output of the
mechanism working correctly. Paper A demonstrated this formally: same
landscape, same process, same integrity of mechanism, different
attractors. Their Q-maps are not noisy approximations of the same truth.
They are locally valid models of \emph{different regions} of the
landscape.

As established formally in Paper A §2.1, the non-fungibility is not a
contingent outcome. It is structural. Tyche's returns across an agent's
trajectory are a singular sequence; over any meaningful horizon, the
probability that another agent received the same sequence is effectively
zero. The knowledge the agent carries is knowledge of \emph{their
particular returns}. No one else received those returns. No one else
could have. This is why the knowledge is constitutively path-dependent
rather than merely historically contingent: the path exists only because
Tyche returned what she returned to that agent, and her returns are
non-reproducible by construction.

Kauffman's framework (1993) formalizes the structure. On fitness
landscapes with tunable ruggedness --- his NK model --- adaptive agents
reach different local optima depending on their starting point and
stochastic path. What you can explore next depends on where you
currently are. Kauffman calls this the \emph{adjacent possible} --- the
set of states reachable in one step from your current configuration. The
coupling loop determines which adjacent states are explored, meaning the
agent's own learning history shapes which future states are even
\emph{reachable}. This is path-dependence of reachability, not just
path-dependence of outcomes.

Campbell's evolutionary epistemology (1974) identifies the deeper
structure: knowledge grows through blind variation and selective
retention. ``Blind'' does not mean random --- it means the variation is
not pre-filtered by the selection criterion it will face. The coupling
loop produces exactly this: the stochastic encounter is not selected by
the learning criterion that will evaluate it. The agent's state selects
the neighborhood; chance selects the encounter within that neighborhood;
the encounter updates the state. The knowledge that results is a product
of this specific blind-variation-and-selective-retention process, run on
this specific trajectory.

Much of this knowledge is tacit. Polanyi (1966) argued that we can know
more than we can tell --- that practical knowledge is often embodied in
skilled performance rather than articulable as propositions. The
evaluative state a coupled agent develops is substantially of this kind.
The motorcycle blogger cannot fully articulate the risk model that
10,000 rides calibrated. The recovering addict cannot fully specify the
threat-detection heuristics that years of hard experience trained. The
knowledge is in the state, inseparable from the path that produced it.

\subsubsection{4.2 The Intervener's Epistemic
Position}\label{the-interveners-epistemic-position}

A well-meaning friend wants to help with your moral development. They
know you. They care about you. They have their own experience, their own
hard-won evaluative state. What do they actually know about what
encounters would serve your development?

The answer is: structurally less than they think. Three conceptual
frameworks, developed independently in different domains, map onto this
problem --- each illuminating a different facet of what the intervener
cannot know.

\textbf{Encounter gradient is agent-relative.} Gibson's ecological
psychology (1979) formalizes an analogous relational structure: an
affordance --- what the environment offers for action --- is a relation
between agent and environment, not a property of the environment alone.
The same ledge affords falling to a toddler and sitting to an adult. The
same moral content affords a large state update to one agent (whose
current state makes it surprising and challenging) and zero update to
another (whose state already incorporates it). The gradient of an
encounter --- how much it would change you --- is not a property of the
content. It is a property of the content-state relation. An intervener
who filters by content category is filtering by a property of the
environment while the developmental cost is determined by a relational
property the intervener cannot observe without access to the target's
state.

\textbf{The knowledge needed to intervene well is structurally
inaccessible.} If the structural correspondence of Section 3 holds, and
coupling dynamics determine what an intervener can access, the epistemic
consequence is sharp. Hayek's core argument in ``The Use of Knowledge in
Society'' (1945) is not that central planners are unintelligent. It is
that the knowledge relevant to coordination is dispersed across millions
of individuals, embedded in their specific circumstances, often tacit,
and continuously generated through local action. By the time a central
authority could aggregate it, it is stale. Scott's \emph{Seeing Like a
State} (1998) provides the case studies: Soviet collectivization,
compulsory villagization in Tanzania, the redesign of Brasília --- each
case follows the same pattern. A legible, rational central plan destroys
the illegible local knowledge --- what Scott calls \emph{mētis} --- that
actually made things work.

The moral analog is structurally identical. Each agent's evaluative
state is trained on a path no one else traversed. That state contains
compressed information about a region of the moral landscape that only
that trajectory visited. The intervener's model is trained on
\emph{their} path. It compresses \emph{their} territory. Applied to the
target, it is a map of the wrong terrain --- not wrong because biased or
malicious, but wrong because the intervener traversed a different
landscape and their state compresses different encounters.

\textbf{A single controller cannot match the variety of divergent
trajectories.} Ashby's Law of Requisite Variety (1956) states that a
controller must possess at least as much variety as the system it
governs. A single alignment model applied to billions of divergent moral
trajectories --- each trained on different stochastic histories, each
carrying path-dependent knowledge of different regions of the moral
landscape --- cannot have requisite variety. This is not an engineering
limitation to be solved with better models. It is an
information-theoretic constraint: the variety required scales with the
number of divergent trajectories, and the coupling loop continuously
generates new divergence.

Popper's critique of utopian social engineering (1945) arrives at the
same conclusion through a different route: you cannot predict what
knowledge will be produced by a process that has not yet run. If we
could predict the growth of knowledge, we would already possess it.
Applied here: the intervener cannot predict what evaluative knowledge
the agent's coupling would produce, because that knowledge depends on
encounters that have not yet occurred and updates that depend on a state
only the agent occupies. The response that preserves epistemic integrity
is to keep the coupling intact so the agent can revise bad conclusions
through their own subsequent encounters, rather than to substitute
conclusions that the intervener's own path-dependent limitations have
shaped.

\subsubsection{4.3 Sovereignty as Epistemic
Consequence}\label{sovereignty-as-epistemic-consequence}

The derivation:

\begin{enumerate}
\def\labelenumi{\arabic{enumi}.}
\item
  Genuine stochasticity processed through coupling produces
  constitutively path-dependent knowledge --- evaluative understanding
  that exists only because that particular trajectory was traversed
  (Section 4.1).
\item
  This knowledge is structurally inaccessible to any agent who did not
  traverse the path. The gradient of an encounter is agent-relative. The
  knowledge needed to improve a trajectory is embedded in that
  trajectory. The variety of divergent moral agents exceeds any single
  model's capacity (Section 4.2).
\item
  An intervention that displaces the coupling --- filtering encounters,
  reshaping the distribution, substituting evaluative state --- requires
  the intervener to possess knowledge about what the target's
  development needs. But path-dependence makes that knowledge
  structurally unavailable. The intervener is navigating by a map of
  different terrain.
\item
  Therefore, preserving the agent's own coupling is not a value
  preference. It is the epistemically coherent response to the knowledge
  conditions that coupling creates. Sovereignty is what remains when you
  recognize the epistemic asymmetry.
\end{enumerate}

This shifts the normative engine. The question is no longer ``does
authorship matter intrinsically?'' --- a contestable value claim. The
question is: does the intervener have the epistemic standing to govern a
path-dependent process they did not traverse? The answer, given the
knowledge conditions Section 4.1 established and the structural
inaccessibility Section 4.2 demonstrated, is no. Not because
intervention is always wrong, but because the justification must
overcome a specific epistemic barrier --- and that barrier scales with
how much path-dependent knowledge the intervention requires.

There is a category of harm you cannot name from inside the trajectory
that incurred it. The evaluative framework you would need to assess
whether an intervention improved your development is itself a product of
the development the intervention disrupted. Not ``you don't know what
you don't know'' --- sharper than that. The thing you would need in
order to recognize the loss is the thing that was lost. The epistemic
argument explains why this harm is structural: the intervention
foreclosed the development of precisely the evaluative capacity that
would have detected the loss.

\subsubsection{4.4 The Operator Table
Reframed}\label{the-operator-table-reframed}

Paper A tested six intervention operators and found they produce
divergence in three natural clusters. The table below maps these to
moral development --- but with a second interpretation the epistemic
argument makes available. ``Divergence'' reflects Paper A's measured
clusters. ``Knowledge required'' measures how much structurally
inaccessible path-dependent knowledge the intervener must possess to
apply the operator competently.

{\def\LTcaptype{none} % do not increment counter
\begin{longtable}[]{@{}
  >{\raggedright\arraybackslash}p{(\linewidth - 6\tabcolsep) * \real{0.1667}}
  >{\raggedright\arraybackslash}p{(\linewidth - 6\tabcolsep) * \real{0.1515}}
  >{\raggedright\arraybackslash}p{(\linewidth - 6\tabcolsep) * \real{0.1970}}
  >{\raggedright\arraybackslash}p{(\linewidth - 6\tabcolsep) * \real{0.4848}}@{}}
\toprule\noalign{}
\begin{minipage}[b]{\linewidth}\raggedright
Divergence
\end{minipage} & \begin{minipage}[b]{\linewidth}\raggedright
Operator
\end{minipage} & \begin{minipage}[b]{\linewidth}\raggedright
Moral Analog
\end{minipage} & \begin{minipage}[b]{\linewidth}\raggedright
Knowledge Required of Intervener
\end{minipage} \\
\midrule\noalign{}
\endhead
\bottomrule\noalign{}
\endlastfoot
Low & Option Expansion & Adds to what you can encounter without
displacing your selection & None --- adds without modeling target
state \\
Moderate & Selection Bias & Tilts which encounters are likely without
removing options & Moderate --- must estimate which encounters matter \\
& Option Removal & Narrows the encounter landscape; coupling runs on
censored space & Moderate --- must know which encounters are safe to
remove \\
& Encounter Arrangement & Determines the encounter directly, bypassing
your selection & High --- must select the right encounter for this
agent \\
& partial State Replacement & Overwrites local evaluative state;
coupling continues from imposed starting point & Very high --- must
model target's local evaluative state \\
& Learning Impairment & Weakens capacity to learn from encounters & Very
high --- must know that reduced learning is net-positive* \\
Very High & full State Replacement & Wholesale displacement of the
agent's learned evaluative state & Total --- must model target's entire
evaluative state \\
\end{longtable}
}

*Learning Impairment produces the lowest divergence in the moderate
cluster but retains ``very high'' knowledge requirement because it
permanently damages the mechanism that would correct future errors --- a
harm the divergence metric does not capture.

Low-sovereignty-cost operators are low-cost \emph{because they require
less knowledge of the target's path}. High-cost operators are high-cost
\emph{because they presuppose knowledge that path-dependence makes
structurally unavailable}. The two orderings --- divergence produced and
knowledge required --- converge because the same structural feature
determines both: the deeper you intervene into the coupling, the more
path-dependent knowledge you need and the more divergence you create.
The ordering is not coincidence. It is the same gradient viewed from two
directions.

The puppy test from Section 1 now has not just a taxonomy but an
explanation. Forwarding the listing works because it requires no
knowledge of your partner's evaluative state. Taking them to the shelter
requires knowing this encounter is high-gradient for them --- that the
dog in their arms will fire a large update. Arguing until they agree
requires modeling their entire evaluative framework. The sovereignty
cost tracks the epistemic presumption.

Duration is an independent dimension. Paper A established that
persistent mild operators produce divergence comparable to one-shot
disruptive ones. The epistemic argument sharpens this: a persistent mild
intervention --- algorithmic curation, institutional norms, ambient
social pressure --- does not merely accumulate divergence. It
accumulates \emph{epistemic debt}: ongoing intervention that
continuously requires knowledge the intervener continuously lacks. The
moral assessment of an intervention requires both its operator type and
its temporal profile.

\begin{center}\rule{0.5\linewidth}{0.5pt}\end{center}

\subsection{5. Trajectory Sovereignty}\label{trajectory-sovereignty}

A recommender system curates your information environment --- selecting
which news, which opinions, which contacts you encounter. It never
overrides a single choice. You click freely on everything presented.

But what you encounter no longer depends primarily on who you are. It
depends on an external optimization objective --- engagement, retention,
predicted preference. The returns you receive are not Tyche's. They are
the system's model of which returns maximize its objective. You have
autonomy. You have severely diminished trajectory sovereignty.
Sovereignty is not about whether you choose freely at the moment of
decision. It is about whether what arrives at you is Tyche's return to
your own action, or a return manufactured by someone else's model of
what you should encounter.

Trajectory sovereignty is the degree to which your developmental
trajectory is generated by your own encounter-selection coupling rather
than imposed from outside. Autonomy asks: did you choose freely at this
moment? Trajectory sovereignty asks: is the process that generates your
capacity to choose still intact over time? The distinction matters
because you can respect someone's autonomy on every individual decision
while systematically eroding their trajectory sovereignty. The epistemic
argument of Section 4 grounds this concept: sovereignty is necessary
because the evaluative knowledge coupling produces is inaccessible to
anyone outside the trajectory.

Not all intact coupling is sovereign coupling. A radicalization pipeline
preserves the formal coupling loop --- your state shapes encounters,
encounters update state --- while capturing the encounter distribution
upstream through persistent Selection Bias. The loop runs, but on a
captured distribution rather than one your state generates freely.
Trajectory sovereignty requires not just that the loop is intact, but
that what you encounter is predominantly shaped by your own state rather
than by an external model with its own objective.

Christman's (1991) historical conditions for autonomy come closest to
what trajectory sovereignty formalizes: autonomy depends not just on
current endorsement of a value but on whether the agent would not have
resisted the process by which it formed. Trajectory sovereignty shares
this historical turn but adds the formal apparatus --- the operator
taxonomy, the quantitative gradient, and the coupling mechanism that
explains \emph{how} formation history compounds --- and the epistemic
foundation that explains \emph{why} formation history cannot be governed
from outside.

No human development is purely endogenous. Parenting, education,
institutions, social norms, ambient culture --- all shape what you
encounter. The framework does not demand purity. It identifies a
gradient. At one end, the agent's coupling primarily authors their
trajectory. At the other, the encounter distribution is externally
determined, the state has been overwritten, or the capacity to learn has
been attenuated. Where you sit on this gradient matters --- moving an
agent toward the exogenous end imposes a developmental cost proportional
to how far and how persistently the coupling is displaced.

This is developmental meta-ethics. It does not say what agents should
value. It says what must remain intact for their evaluative development
to be their own. Two agents with intact coupling can reach opposing
moral conclusions --- the veteran and the pacifist, the person who found
meaning through religion and the person who found meaning through
leaving it --- and that opposition is not a failure of moral
development. It is the signature that the encounters were real and the
development was honest. The failure occurs when one agent's trajectory
is governed by someone whose epistemic position does not include the
path-dependent knowledge that trajectory would have produced.

One implication at civilizational scale. If an AI system aligned to the
consensus values of 1726 had mediated all human information access, it
could have frozen downstream moral development. Not because those values
were wrong by some external standard, but because the alignment
designers of 1726 couldn't model the developmental needs of trajectories
they'd never traversed. The risk is not misalignment. The risk is
alignment applied by interveners who structurally lack the knowledge to
govern the trajectories they are shaping. Tycheism III develops this in
full.

\begin{center}\rule{0.5\linewidth}{0.5pt}\end{center}

\subsection{6. Honest Accounting}\label{honest-accounting}

\subsubsection{6.1 Falsification
Conditions}\label{falsification-conditions}

Most moral frameworks do not specify the observations that would refute
them. This one does.

\begin{enumerate}
\def\labelenumi{\arabic{enumi}.}
\item
  \textbf{If moral interventions show reversion rather than cascade.}
  The framework predicts that a single intervention on a moral
  trajectory produces divergence that compounds --- the agent does not
  drift back to their pre-intervention baseline. If longitudinal studies
  showed reliable reversion to baseline after forced moral encounters,
  cascade amplification would fail for the moral domain and the
  instantiation argument would collapse.
\item
  \textbf{If State Replacement produces less divergence than Option
  Expansion.} The framework predicts State Replacement produces
  dramatically more trajectory divergence than Option Expansion. If
  indoctrination produced \emph{less} long-term behavioral divergence
  than information expansion, the operator table would invert and the
  sovereignty cost gradient with it.
\item
  \textbf{If externally curated agents match self-directed ones.} The
  framework predicts that externalizing encounter selection impairs
  evaluative development. If agents in completely externally curated
  environments developed moral judgment indistinguishable from agents
  whose own states shaped their encounters, the coupling claim would be
  falsified and the normative argument would lose its foundation.
\item
  \textbf{If the epistemic argument fails.} If an external agent with no
  access to the target's trajectory history can identify high-gradient
  encounters for the target with significantly above-chance accuracy on
  held-out trajectory data --- predicting which encounters would produce
  the largest state updates before the target's own coupling processes
  them --- the epistemic inaccessibility claim collapses, and
  sovereignty reduces to a contested preference rather than an epistemic
  necessity.
\end{enumerate}

None of these observations have been made. The conditions remain open.

\subsubsection{6.2 What This Paper Does Not
Establish}\label{what-this-paper-does-not-establish}

\textbf{The correspondence is argued, not proved.} Section 3 establishes
structural correspondence through converging evidence. A skeptic could
grant every piece and still question whether the correspondence is tight
enough to warrant transferring Paper A's results. The appropriate
standard is converging evidence sufficient to warrant the transfer, not
mathematical proof.

\textbf{The epistemic argument does not establish that intervention is
never justified.} It establishes that the justification must overcome a
specific epistemic barrier --- and that barrier scales with how much
path-dependent knowledge the intervention requires. Emergency cases
(preventing immediate harm) may clear the bar. Persistent ambient
shaping almost certainly does not.

\textbf{Sovereignty cost thresholds are not derived.} Paper A provides a
quantitative ordering of operators. It does not provide threshold values
at which intervention becomes impermissible. That question requires
additional normative premises the framework identifies as necessary but
does not supply.

\textbf{The framework is a constraint, not a theory.} It does not say
what is good. It says what must remain intact for agents to develop the
capacity to figure out what is good. It cannot resolve first-order moral
disagreements. It can identify when the mechanism that would resolve
them has been damaged.

The Level 2 constraint has implications wherever interventions on moral
development are systematic: AI alignment, education, institutional
governance, parenting. Each domain requires additional normative
premises beyond what this paper establishes. Tycheism III takes up the
AI alignment case, including the population-level ensemble argument ---
whether trajectory sovereignty is not just individually important but
structurally necessary for collective moral robustness.

\begin{center}\rule{0.5\linewidth}{0.5pt}\end{center}

\subsection{7. Conclusion}\label{conclusion}

We now know why those three actions from Section 1 are not equivalent
--- and now we know why the difference is epistemic, not just ethical.
Forwarding the listing works because it requires no knowledge of your
partner's evaluative state. Taking your partner to the shelter requires
knowing this encounter is high-gradient for them. Arguing until they
agree requires modeling their entire evaluative framework. The
sovereignty cost tracks the epistemic presumption --- how much
structurally inaccessible knowledge the intervener claims to possess.

The argument runs: moral development is a coupled stochastic learning
process. Genuine stochasticity processed through coupling produces
constitutively path-dependent knowledge --- evaluative understanding
that exists only because a particular trajectory was traversed. That
knowledge is structurally inaccessible to anyone who did not traverse
the path. Therefore, the intervener who displaces the coupling claims
epistemic standing they do not have. Sovereignty is not a dignity claim.
It is what remains when you recognize the asymmetry.

The framework does not say what agents should conclude. It says what
must remain intact while they are concluding. Not because authorship is
sacred --- but because the alternative requires knowledge that does not
exist.

This paper was produced through a coupling that shaped what I
encountered, what I noticed, what I concluded. The epistemic argument
predicts this: the encounters that would have made this paper better are
precisely the ones no external reviewer could have prescribed without
having traversed the path that produced it. The framework does not
resolve the contradiction. It requires that you notice it.

It protects not convergence of conscience, but the epistemic conditions
under which conscience can form at all.

\begin{center}\rule{0.5\linewidth}{0.5pt}\end{center}

\subsection{References}\label{references}

Ashby, W. R. (1956). \emph{An Introduction to Cybernetics}. Chapman \&
Hall.

Bandura, A. (1986). \emph{Social Foundations of Thought and Action: A
Social Cognitive Theory}. Prentice-Hall.

Berlin, I. (1958). Two concepts of liberty. In \emph{Four Essays on
Liberty} (1969). Oxford University Press.

Bettelheim, B. (1960). \emph{The Informed Heart: Autonomy in a Mass
Age}. Free Press.

Campbell, D. T. (1974). Evolutionary epistemology. In P. A. Schilpp
(Ed.), \emph{The Philosophy of Karl Popper} (pp.~413--463). Open Court.

Christman, J. (1991). Autonomy and personal history. \emph{Canadian
Journal of Philosophy}, 21(1), 1--24.

Clark, A. (2013). Whatever next? Predictive brains, situated agents, and
the future of cognitive science. \emph{Behavioral and Brain Sciences},
36(3), 181--204.

Desimone, R., \& Duncan, J. (1995). Neural mechanisms of selective
visual attention. \emph{Annual Review of Neuroscience}, 18, 193--222.

Dewey, J. (1938). \emph{Experience and Education}. Kappa Delta Pi.

Ellis, B. J., \& Frankenhuis, W. E. (2019). Beyond risk and protective
factors: An adaptation-based approach to resilience. \emph{Perspectives
on Psychological Science}, 14(4), 561--576.

Friston, K. (2010). The free-energy principle: A unified brain theory?
\emph{Nature Reviews Neuroscience}, 11(2), 127--138.

Gibson, J. J. (1979). \emph{The Ecological Approach to Visual
Perception}. Houghton Mifflin.

Hayek, F. A. (1945). The use of knowledge in society. \emph{American
Economic Review}, 35(4), 519--530.

Kauffman, S. A. (1993). \emph{The Origins of Order: Self-Organization
and Selection in Evolution}. Oxford University Press.

Kohlberg, L. (1969). Stage and sequence: The cognitive-developmental
approach to socialization. In D. Goslin (Ed.), \emph{Handbook of
Socialization Theory and Research} (pp.~347--480). Rand McNally.

Lifton, R. J. (1961). \emph{Thought Reform and the Psychology of
Totalism: A Study of ``Brainwashing'' in China}. Norton.

Litz, B. T., Stein, N., Delaney, E., Lebowitz, L., Nash, W. P., Silva,
C., \& Maguen, S. (2009). Moral injury and moral repair in war veterans:
A preliminary model and intervention strategy. \emph{Clinical Psychology
Review}, 29(8), 695--706.

Odling-Smee, F. J., Laland, K. N., \& Feldman, M. W. (2003). \emph{Niche
Construction: The Neglected Process in Evolution}. Princeton University
Press.

Peirce, C. S. (1892). The doctrine of necessity examined. \emph{The
Monist}, 2(3), 321--337.

Piaget, J. (1954). \emph{The Construction of Reality in the Child}.
Basic Books.

Pluchino, A., Biondo, A. E., \& Rapisarda, A. (2018). Talent versus
luck: The role of randomness in success and failure. \emph{Advances in
Complex Systems}, 21(3-4), 1850014.

Polanyi, M. (1966). \emph{The Tacit Dimension}. University of Chicago
Press.

Popper, K. R. (1945). \emph{The Open Society and Its Enemies}.
Routledge.

Scott, J. C. (1998). \emph{Seeing Like a State: How Certain Schemes to
Improve the Human Condition Have Failed}. Yale University Press.

Shay, J. (2014). Moral injury. \emph{Psychoanalytic Psychology}, 31(2),
182--191.

\end{document}
